% GENERATED FILE. Edit canonical lessons, then run npm run generate:books.
% canonical-source-hash: fnv1a64:b6e09dba77776fc3
% canonical-lessons: JA-C409-chiiki, JA-C409-shuto, JA-C409-kokkyou, JA-C409-sakai, JA-C409-hougen

\chapter{Regions and Borders}
\label{ch:ja-regions-and-borders}
\hlchaptermodality{409}

\begin{chapteropening}
\textbf{By the end of this chapter:} \emph{I can talk about regions, borders and dialects.}

Complete the last lesson of chapter 409: I can talk about regions, borders and dialects.
\end{chapteropening}

\section[\texorpdfstring{\ja{ちいき} (chiiki)}{ちいき (chiiki)}]{\ja{ちいき} (chiiki) — an area; a district}

\label{lesson:JA-C409-chiiki}



\begin{quote}
Before the new one: say the Japanese for religion, then the Japanese for to enshrine.
\end{quote}

\subsection*{You'll want to know: \ja{ちいき}}
\textbf{\ja{ちいき}} — \emph{chiiki} — \textquotedblleft{}an area; a district\textquotedblright{}.

\begin{quote}
\textbf{\ja{この} \ja{ちいきは} \ja{ゆきが} \ja{おおいです。}} — \emph{kono chiiki wa yuki ga ōi desu} — \textquotedblleft{}This region gets a lot of snow.\textquotedblright{}
\end{quote}

\begin{cousinweb}
\textbf{\ja{ちいき}} (\emph{chiiki}) is Sino-Japanese (kango): two Chinese characters read with their on'yomi, readings borrowed from Chinese. They are \emph{chi}, \textquotedblleft{}ground\textquotedblright{} (Middle Chinese \emph{dijH}) and \emph{iki}, \textquotedblleft{}region\textquotedblright{} (\emph{hwik}) (KANJIDIC2; Middle Chinese in Baxter's transcription).
\end{cousinweb}

\begin{grammarlens}[title={What you've built}]
One more word for getting around.
\end{grammarlens}

\subsection*{Your turn}
\begin{itemize}
\raggedright
  \item \emph{Say it:} \emph{chiiki}
  \item \emph{Say it:} \emph{chiiki}, once more
  \item \emph{Say it:} \emph{kono chiiki wa yuki ga ōi desu} — This region gets a lot of snow.
  \item \emph{Recall:} say the Japanese for religion, then the Japanese for to enshrine, then say \emph{chiiki} again
\end{itemize}

\begin{tcolorbox}[breakable,colback=teal!4,colframe=teal!35!black,title={Before you move on}]
What does \ja{ちいき} mean? (\textquotedblleft{}An area; a district\textquotedblright{}.) Say it once more.
\end{tcolorbox}
\section[\texorpdfstring{\ja{しゅと} (shuto)}{しゅと (shuto)}]{\ja{しゅと} (shuto) — a capital city}

\label{lesson:JA-C409-shuto}



\begin{quote}
Before the new one: say the Japanese for to enshrine, then the Japanese for an area.
\end{quote}

\subsection*{You'll want to know: \ja{しゅと}}
\textbf{\ja{しゅと}} — \emph{shuto} — \textquotedblleft{}a capital city\textquotedblright{}.

\begin{quote}
\textbf{\ja{しゅとの} \ja{ちかてつは} \ja{いつも} \ja{こんざつしています。}} — \emph{shuto no chikatetsu wa itsumo konzatsu shite imasu} — \textquotedblleft{}The subway in the capital is always crowded.\textquotedblright{}
\end{quote}

\begin{cousinweb}
\textbf{\ja{しゅと}} (\emph{shuto}) is Sino-Japanese (kango): two Chinese characters read with their on'yomi, readings borrowed from Chinese. They are \emph{shu}, \textquotedblleft{}neck\textquotedblright{} (Middle Chinese \emph{syuwX}) and \emph{to}, \textquotedblleft{}capital\textquotedblright{} (\emph{tu}) (KANJIDIC2; Middle Chinese in Baxter's transcription).
\end{cousinweb}

\begin{grammarlens}[title={What you've built}]
One more word for getting around.
\end{grammarlens}

\subsection*{Your turn}
\begin{itemize}
\raggedright
  \item \emph{Say it:} \emph{shuto}
  \item \emph{Say it:} \emph{shuto}, once more
  \item \emph{Say it:} \emph{shuto no chikatetsu wa itsumo konzatsu shite imasu} — The subway in the capital is always crowded.
  \item \emph{Recall:} say the Japanese for to enshrine, then the Japanese for an area, then say \emph{shuto} again
\end{itemize}

\begin{tcolorbox}[breakable,colback=teal!4,colframe=teal!35!black,title={Before you move on}]
What does \ja{しゅと} mean? (\textquotedblleft{}A capital city\textquotedblright{}.) Say it once more.
\end{tcolorbox}
\section[\texorpdfstring{\ja{こっきょう} (kokkyō)}{こっきょう (kokkyō)}]{\ja{こっきょう} (kokkyō) — a national border}

\label{lesson:JA-C409-kokkyou}



\begin{quote}
Before the new one: say the Japanese for an area, then the Japanese for a capital city.
\end{quote}

\subsection*{You'll want to know: \ja{こっきょう}}
\textbf{\ja{こっきょう}} — \emph{kokkyō} — \textquotedblleft{}a national border\textquotedblright{}.

\begin{quote}
\textbf{\ja{くるまで} \ja{こっきょうを} \ja{こえました。}} — \emph{kuruma de kokkyō o koemashita} — \textquotedblleft{}We crossed the border by car.\textquotedblright{}
\end{quote}

\begin{cousinweb}
\textbf{\ja{こっきょう}} (\emph{kokkyō}) is Sino-Japanese (kango): two Chinese characters read with their on'yomi, readings borrowed from Chinese. They are \emph{koku}, \textquotedblleft{}country\textquotedblright{} (Middle Chinese \emph{kwok}) and \emph{kyō}, \textquotedblleft{}boundary\textquotedblright{} (\emph{kjaengX}), joined as \emph{kokkyō} (KANJIDIC2; Middle Chinese in Baxter's transcription).
\end{cousinweb}

\begin{grammarlens}[title={What you've built}]
One more word for getting around.
\end{grammarlens}

\subsection*{Your turn}
\begin{itemize}
\raggedright
  \item \emph{Say it:} \emph{kokkyō}
  \item \emph{Say it:} \emph{kokkyō}, once more
  \item \emph{Say it:} \emph{kuruma de kokkyō o koemashita} — We crossed the border by car.
  \item \emph{Recall:} say the Japanese for an area, then the Japanese for a capital city, then say \emph{kokkyō} again
\end{itemize}

\begin{tcolorbox}[breakable,colback=teal!4,colframe=teal!35!black,title={Before you move on}]
What does \ja{こっきょう} mean? (\textquotedblleft{}A national border\textquotedblright{}.) Say it once more.
\end{tcolorbox}
\section[\texorpdfstring{\ja{さかい} (sakai)}{さかい (sakai)}]{\ja{さかい} (sakai) — a boundary}

\label{lesson:JA-C409-sakai}



\begin{quote}
Before the new one: say the Japanese for a capital city, then the Japanese for a national border.
\end{quote}

\subsection*{You'll want to know: \ja{さかい}}
\textbf{\ja{さかい}} — \emph{sakai} — \textquotedblleft{}a boundary\textquotedblright{}.

\begin{quote}
\textbf{\ja{かわが} \ja{ふたつの} \ja{まちの} \ja{さかいです。}} — \emph{kawa ga futatsu no machi no sakai desu} — \textquotedblleft{}The river is the boundary between the two towns.\textquotedblright{}
\end{quote}

\begin{cousinweb}
\textbf{\ja{さかい}} (\emph{sakai}) is a native Japanese word (wago), inherited rather than borrowed from Chinese; its kanji was attached later as a kun'yomi (JMdict).
\end{cousinweb}

\begin{grammarlens}[title={What you've built}]
One more word for getting around.
\end{grammarlens}

\subsection*{Your turn}
\begin{itemize}
\raggedright
  \item \emph{Say it:} \emph{sakai}
  \item \emph{Say it:} \emph{sakai}, once more
  \item \emph{Say it:} \emph{kawa ga futatsu no machi no sakai desu} — The river is the boundary between the two towns.
  \item \emph{Recall:} say the Japanese for a capital city, then the Japanese for a national border, then say \emph{sakai} again
\end{itemize}

\begin{tcolorbox}[breakable,colback=teal!4,colframe=teal!35!black,title={Before you move on}]
What does \ja{さかい} mean? (\textquotedblleft{}A boundary\textquotedblright{}.) Say it once more.
\end{tcolorbox}
\section[\texorpdfstring{\ja{ほうげん} (hōgen)}{ほうげん (hōgen)}]{\ja{ほうげん} (hōgen) — a dialect}

\label{lesson:JA-C409-hougen}



\begin{quote}
Before the new one: say the Japanese for a national border, then the Japanese for a boundary.
\end{quote}

\subsection*{You'll want to know: \ja{ほうげん}}
\textbf{\ja{ほうげん}} — \emph{hōgen} — \textquotedblleft{}a dialect\textquotedblright{}.

\begin{quote}
\textbf{\ja{そふは} \ja{いなかの} \ja{ほうげんで} \ja{はなします。}} — \emph{sofu wa inaka no hōgen de hanashimasu} — \textquotedblleft{}My grandfather speaks a country dialect.\textquotedblright{}
\end{quote}

\begin{cousinweb}
\textbf{\ja{ほうげん}} (\emph{hōgen}) is Sino-Japanese (kango): two Chinese characters read with their on'yomi, readings borrowed from Chinese. They are \emph{hō}, \textquotedblleft{}direction\textquotedblright{} (Middle Chinese \emph{pjang}) and \emph{gen}, \textquotedblleft{}say\textquotedblright{} (\emph{ngjon}) (KANJIDIC2; Middle Chinese in Baxter's transcription).
\end{cousinweb}

\begin{grammarlens}[title={What you've built}]
One more word for getting around.
\end{grammarlens}

\subsection*{Your turn}
\begin{itemize}
\raggedright
  \item \emph{Say it:} \emph{hōgen}
  \item \emph{Say it:} \emph{hōgen}, once more
  \item \emph{Say it:} \emph{sofu wa inaka no hōgen de hanashimasu} — My grandfather speaks a country dialect.
  \item \emph{Recall:} say the Japanese for a national border, then the Japanese for a boundary, then say \emph{hōgen} again
\end{itemize}

\begin{tcolorbox}[breakable,colback=teal!4,colframe=teal!35!black,title={Before you move on}]
What does \ja{ほうげん} mean? (\textquotedblleft{}A dialect\textquotedblright{}.) Say it once more.
\end{tcolorbox}
